\documentclass[11pt]{article}
\usepackage[margin=1in]{geometry}
\usepackage{amsmath,amssymb,amsthm}
\usepackage{booktabs}
\usepackage{hyperref}
\usepackage{natbib}
\hypersetup{colorlinks=true,linkcolor=blue,citecolor=blue,urlcolor=blue}

\title{Clifford networks do not improve Ramsey sensitivity}

\author{Akshit\thanks{Code: \url{https://github.com/akshitj11/clifford-based-network}.
Correspondence: \texttt{akshitcasual1@gmail.com}.}}

\date{October 2026}

\begin{document}
\maketitle

\begin{abstract}
A Clifford layer does not improve quantum sensitivity.
On a simulated planar Ramsey magnetometer it matches a least-squares inversion of the known likelihood, sits on the Cram\'er--Rao bound at 512 shots ($\mathrm{rms}/\sqrt{\mathrm{CRB}}\approx 0.99$), and loses to a plain linear net on hard-iron / soft-iron calibration.
The one place it is decisive is the product a vector sensor already computes, $(\mathrm{readout}-\tfrac12)\times\mathrm{axis}$.
Given the raw pieces and not the product, a 68-parameter $\mathrm{Cl}(2,0)$ net is at the noise floor with 20 labeled fields (field MSE $\approx 0.002$).
A parameter-matched $\tanh$ MLP is still at $0.32$ after 500 fields.
That is a decoder result.
The qubit does not get quieter.
\end{abstract}

\section{The claim}

Sensing is the extraction of a field $B$ from shots~\citep{degen2017,helstrom1969}.
On a single-fringe Ramsey protocol that extraction is already solved by $\arcsin$ and a $2\times 2$ solve~\citep{ramsey1950,taylor2008}.
A network that matches that solve has inverted the record.
It has not improved the sensor.

Clifford (geometric-algebra) layers are a natural candidate for the decoder anyway.
A magnetic field is a grade-1 object.
A vector magnetometer reports it in the sensor frame, so a platform rotation $R$ sends $B\mapsto RB$ and the measurement axes with it.
The geometric product respects that action: $R(ab)\tilde R=(Ra\tilde R)(Rb\tilde R)$~\citep{doran2003,ruhe2023cgenn}.
Spin-$\tfrac12$ precession is a rotor in $\mathrm{Cl}(3,0)$.
This note stays in $\mathrm{Cl}(2,0)$.
A 3D Bloch magnetometer is a later algebra, not a better first measurement.

Machine-learning inverses for atomic and NV magnetometers already exist~\citep{liu2023,zhang2024}.
Equivariant and physics-informed nets for platform magnetometry exist too.
Clifford networks have been shown on $n$-body, fluids, and Lorentz-equivariant HEP~\citep{ruhe2023gcan,ruhe2023cgenn,brandstetter2023}.
What is thin is the combination: a Clifford layer, a Ramsey likelihood, a classical estimator, and a Cram\'er--Rao number on the same simulated shots.
That is the measurement this note reports.

The first sweep is a tie by construction.
Once any model is handed the vector observation $v_k=(P_k-\tfrac12)u_k$, every estimator is doing the same linear inversion.
The architectures only come apart when that product is taken away, or when the classical fit believes the wrong sensor.

\section{Protocol}

A spin-$\tfrac12$ runs a Ramsey sequence along each of $m=6$ unit axes $u_k$, biased to the steepest point of the fringe (control phase $-\pi/2$):
\begin{align}
\phi_k &= \gamma T\,(B\cdot u_k),\\
P_k &= \tfrac12\bigl(1+V\sin\phi_k\bigr)+\mathrm{noise},\qquad
V=\exp(-T/T_2).
\end{align}
Default numbers: $\gamma=1$, $T=1$, $T_2=5$ so $V\approx 0.82$, Gaussian readout $\sigma=0.02$ unless replaced by binomial shot noise.
Fields are drawn uniformly in the disk $|B|\le 0.8$, which keeps $|\phi|\le 0.8\,\mathrm{rad}$.
$\sin$ is monotonic there.
The record is invertible.
That is the regime in which least squares is allowed to be strong~\citep{taylor2008,maze2008}.

An optional device defect is applied before projection,
$B_{\mathrm{eff}}=\mathrm{diag}(s)B+h$,
with soft-iron gain $s$ (default $(1,1)$) and hard-iron bias $h$ (default $0$).
This is the usual hard-iron / soft-iron (Tolles--Lawson) distortion of a vector magnetometer.
Defects are a property of the sensor, constant across samples.
Fields and orientations vary.

The classical estimator inverts the fringe,
\begin{equation}
\mathrm{proj}_k=\arcsin\bigl(\mathrm{clip}((2P_k-1)/V)\bigr)/(\gamma T),
\end{equation}
and solves the $2\times 2$ normal equations for $B_{\mathrm{eff}}$.
An oracle then undoes the true $(s,h)$.
A naive fit assumes $s=1$, $h=0$, and the nominal $T_2$.

Learned models share one topology: a shared per-axis message, a sum over the six axes, a shared update that reads a 2D field.
\texttt{CliffordMag} carries $\mathrm{Cl}(2,0)$ multivector channels through geometric products.
\texttt{FlatMag} is a parameter-matched real linear net ($68$ vs $70$ scalars).
Training is full-batch analytic SGD, $400$ epochs, learning rate $0.1$, three seeds unless noted.
The stack is CPU-only Rust, single-threaded, hand-written geometric-product VJPs, splitmix64 data.
A sweep finishes in minutes.

Field MSE is the mean of $(\Delta B_x)^2+(\Delta B_y)^2$ on $500$ held-out samples.
RMS field error is its square root.
Equivariance error is the RMS gap between rotating the estimate and estimating on the rotated apparatus.

\section{Clean inversion is a tie}

On the default Gaussian protocol, least squares, Clifford, and flat all report field MSE $\approx 0.00191$ on the held-out set.
At $1000$ training fields the two nets are $0.001877$ (mean over seeds).
That is the readout-noise floor of this protocol, not a model ranking.

The slope of the fringe near zero field is $dP/dB\approx\tfrac12 V\approx 0.41$.
$\sigma=0.02$ is then about $0.05$ of field error on one axis.
Six axes, least-squares combined, put the two-component MSE near $0.0016$ in the linear approximation.
The measured $0.00191$ is a little worse because $\sin$ flattens at larger phase and $T_2$ has already cut $18\%$ of the contrast.

What splits the nets is rotation, not accuracy.
At $N=1000$, equivariance error is $0.00092$ (Clifford) versus $0.00152$ (flat).
At $N=100$ the same gap is $0.0044$ versus $0.0067$.
The classical inversion is at $6\times 10^{-8}$.
Clifford is the closer of the two learned models to the group action.
It is not a more accurate field estimate.

\section{A wrong classical model is beatable. Geometry is not why.}

A real device has a bias, a gain, and a $T_2$ you do not know exactly.
The honest baseline is a lab fit that assumes an ideal sensor.

Defect level $d$ sets
$h=(0.2d,-0.15d)$,
$s=(1+0.3d,\,1-0.2d)$,
$T_2=5(1-0.2d)$,
and the naive fit still believes $h=0$, $s=1$, $T_2=5$.

\begin{table}[h]
\centering
\caption{Held-out field MSE versus defect level $d$. Oracle least squares is given the true $(s,h,T_2)$.}
\label{tab:misspec}
\begin{tabular}{lcccc}
\toprule
$d$ & naive LS & oracle LS & Clifford & flat \\
\midrule
$0.0$ & $0.00191$ & $0.00191$ & $0.00188$ & $0.00188$ \\
$0.5$ & $0.0221$  & $0.00199$ & $0.00611$ & $0.00207$ \\
$1.0$ & $0.0761$  & $0.00240$ & $0.0171$  & $0.00278$ \\
\bottomrule
\end{tabular}
\end{table}

At $d=1$ the naive fit is about $40\times$ the noise floor.
The oracle is essentially back on the floor.
The flat net tracks the oracle ($0.00278$).
Clifford beats the naive fit ($0.017$ versus $0.076$) and loses to the flat net by about $6\times$.

The defect is a world-frame diagonal scale plus an offset.
That is not a rotor.
A linear map on $v_k$ plus a constant channel can absorb $s$ and $h$.
The flat net does that.
Learned models beat a misspecified likelihood.
The architecture that won is the one without geometric structure.

Equivariance error rises with $d$ for every estimator that tries to undo the defect ($0.30$ for the oracle at $d=1$).
Soft-iron in world axes breaks the rotation symmetry of the record.
The naive fit stays numerically ``equivariant'' because it is consistently wrong in a rotation-covariant way.

\section{Take the product away}

\texttt{CliffordMag} and \texttt{FlatMag} were handed $v_k=(P_k-\tfrac12)u_k$.
A single linear layer can finish the job.
That is why they tied at $N=100$.

\texttt{CliffordMagRaw} is given two separate channels, $\mathrm{scalar}(P_k-\tfrac12)$ and $\mathrm{vector}(u_k)$, and forms the observable with one geometric product.
No weights in that product.
The rest of the net is the same $68$-parameter message/update pair.
\texttt{FlatMlpMag} sees $[P_k-\tfrac12,\,u_x,\,u_y]$ through a $\tanh$ hidden layer ($67$ parameters) and has to learn the bilinear map.
Five seeds, $600$ epochs, training sizes $20/50/100/500$.

\begin{table}[h]
\centering
\caption{Field MSE when the model must form $(P-\tfrac12)u$ itself. Mean $\pm$ std over five seeds.}
\label{tab:efficiency}
\begin{tabular}{lcc}
\toprule
$N$ & Clifford raw & $\tanh$ MLP \\
\midrule
$20$  & $0.00196\pm 0.00004$ & $0.432\pm 0.057$ \\
$50$  & $0.00195\pm 0.00005$ & $0.353\pm 0.022$ \\
$100$ & $0.00192\pm 0.00003$ & $0.332\pm 0.006$ \\
$500$ & $0.00188\pm 0.00001$ & $0.323\pm 0.002$ \\
\bottomrule
\end{tabular}
\end{table}

$0.32$ is the MSE of a model that predicts near-zero field on this disk.
The MLP never forms the product in this budget.
Clifford is on the floor at the first point in the sweep.
Equivariance follows the same split.
At $N=20$, Clifford is $0.004$--$0.010$; the MLP is $0.18$--$0.40$.
At $N=500$, Clifford is $0.0002$--$0.002$; the MLP is still $0.02$--$0.06$.

The geometric product is the map the sensor physics already uses.
A generic nonlinearity of the same width does not find it from a few hundred examples.
There is no crossover in the range we ran.
Saying the gap closes as $N$ grows would be a prediction, not a measurement.
A wider MLP, an explicit multiply, or a kernel would be a stronger baseline.
Those are not in this sweep.

\section{Shot noise against the bound}

Gaussian $\sigma$ is a convenience.
A projective readout of $n$ shots is a binomial count.
Fisher information for this Ramsey model, per axis, is
\begin{align}
\frac{dP}{dB_{\mathrm{proj}}}&=\tfrac12 V\gamma T\cos\phi_k,\\
\mathrm{Var}(\hat P)&=P_k(1-P_k)/n,\\
F&=\sum_k \frac{(dP/dB)(dP/dB)^\top}{\mathrm{Var}_k},\qquad
\mathrm{CRB}=\mathrm{tr}(F^{-1}).
\end{align}
The figure that should travel is $\mathrm{rms}/\sqrt{\mathrm{CRB}}$, plus the shot-scaled sensitivity $\gamma T\sqrt{n}\cdot\mathrm{rms}$~\citep{degen2017,helstrom1969}.
A ratio near $1$ means the estimator used the shots.
A ratio of $0.5$ would be a claim of super-classical information and a reason to stop and find the bug.

\begin{table}[h]
\centering
\caption{RMS field error versus $\sqrt{\mathrm{CRB}}$ under binomial readout.}
\label{tab:shots}
\begin{tabular}{lcccccc}
\toprule
$n$ & class.\ rms & Cliff.\ rms & flat rms & $\sqrt{\mathrm{CRB}}$ & class./bnd & Cliff./bnd \\
\midrule
$32$  & $0.192$  & $0.169$  & $0.169$  & $0.182$  & $1.056$ & $0.929$ \\
$128$ & $0.0925$ & $0.0898$ & $0.0899$ & $0.0910$ & $1.016$ & $0.987$ \\
$512$ & $0.0448$ & $0.0449$ & $0.0449$ & $0.0454$ & $0.986$ & $0.988$ \\
\bottomrule
\end{tabular}
\end{table}

Error scales as $1/\sqrt{n}$.
Sixteen times more shots ($32\to 512$) cut RMS by about four ($0.192\to 0.045$).
Sensitivity sits at $1.01$--$1.09$ in these units.
At $512$ shots all three estimators are on the bound.

At $32$ shots the nets sit about $7\%$ under the unbiased CRB.
The Cram\'er--Rao statement is for unbiased estimators~\citep{helstrom1969}.
A net trained on MSE shrinks toward the mean of the training disk.
That can beat an unbiased bound on squared error without extracting more information than the likelihood contains.
The classical inversion stays just above the bound ($1.056$), which is the shape you want from a correctly derived Fisher matrix.

Nobody here improved the sensor.
They used the shots.

\section{What is allowed}

On a known Ramsey likelihood, a Clifford net is an efficient function class for the inverse, not a better estimator than least squares.
When the inverse requires multiplying a scalar readout onto a direction, that product is cheap in $\mathrm{Cl}(2,0)$ and expensive for a small flat MLP.
When the inverse requires undoing a world-frame gain and bias, a flat net is the better function class.
Under binomial readout, both nets match the Cram\'er--Rao scaling.
Sensitivity is set by $n$, $V$, and $T$, not by the layer type.

The claim that does not survive: that a Clifford network is an upgrade to quantum sensing.

Simulation only.
There is no NV-center, fluxgate, or SQUID dataset in the repository.
Noise and $T_2$ are modeled; they are not measured.
The sensor is an ideal planar vector magnetometer.
An NV-center has fixed crystal axes~\citep{zhang2024}.
Full rotation equivariance does not hold for that device.
You would pass the axes in as data, which this model already does, and you would stop claiming the network is equivariant by construction.

The operating point stays inside one fringe.
Multiple evolution times, phase wrapping, and adaptive Ramsey are out of scope.
Those are the regimes where a learned inverse sometimes earns its keep against a misspecified or multimodal likelihood.
They are also the regimes where a classical estimator has to become a proper maximum-likelihood search, not a $2\times 2$ solve.

$\mathrm{Cl}(3,0)$ is not implemented.
A 3D Bloch magnetometer on the same easy inverse would reproduce the clean-data tie.
It becomes interesting only after the product-removal and shot-noise tests are restated in 3D.

Reproduce with
\texttt{cargo test} and the four release binaries
\texttt{train\_magnetometer},
\texttt{train\_mag\_misspec},
\texttt{train\_mag\_efficiency},
\texttt{train\_mag\_shots}.
CSVs land in \texttt{results/}.

The next measurement that could still change the story is not ``add $\mathrm{Cl}(3,0)$ and rerun the easy inverse.''
It is a wrapped-fringe, multi-$T$ protocol, or a real vector-magnetometer log, with the same three estimators and the same refusal to call a noise-floor tie a sensing win.

\bibliographystyle{plainnat}
\bibliography{refs}

\end{document}
